%!TEX program = lualatex
\documentclass[ngerman,headheight=70pt]{scrartcl}
\usepackage[ngerman]{babel}
\usepackage{blindtext}
\usepackage[utf8]{luainputenc}
\usepackage{scrlayer-scrpage}
\usepackage[T1]{fontenc}
\usepackage{fontspec}
\usepackage{multicol}
\usepackage{enumitem}
\usepackage[german=quotes]{csquotes}
\usepackage[bottom=2cm,footskip=8mm]{geometry}
\usepackage{pdfpages}
\setmainfont{Source Sans Pro}
\MakeOuterQuote{"}

\setlength{\parindent}{0pt}
\setlength{\parskip}{6pt}

\pagestyle{plain.scrheadings}
\thispagestyle{scrheadings}
\lohead{\vspace*{-1.5cm}\includegraphics[scale=0.4]{../images/stupa_logo_rot.pdf}}
\rohead{\\Präsidium des\\Studierendenparlaments}

\lofoot[Universität Hamburg · Präsidium des Studierendenparlaments\\\hrule Von-Melle-Park 5 · D-20146 Hamburg]{Universität Hamburg · Präsidium des Studierendenparlaments\\\hrule Von-Melle-Park 5 · D-20146 Hamburg}
\cofoot[]{}
\cohead[\thepage]{}

\renewcommand{\thesubsection}{TOP \arabic{subsection}}
\renewcommand{\thesubsubsection}{\arabic{subsubsection}.}
\setcounter{subsection}{-1}

\begin{document}
    UHH · StuPa-Präsidium · Von-Melle-Park 5 · D-20146 Hamburg

    %%% PAGE ONE %%%%
    \section*{Protokoll der 1. Fortsetzung der 1. Sitzung des Studierendenparlaments vom 28. April 2016}

    \textbf{Protokoll: Jim Martens}\\
    \textbf{Ort: Raum S08 in der ehemaligen HWP }\\
    \textbf{Beginn: 18.25 Uhr}\\
    \textbf{Ende: 23.41 Uhr}

    \vspace{0.5cm}
    \begin{tabular}{ll}
        Anwesend: & \\
            RCDS (5 Sitze): & Ramon Weilinger, Antonia Niecke, Ramin Shakiba, \\
                            & Jennifer Maack, Benjamin Welling \\
             CampusGrün (14 Sitze): & Laura Franzen, Geoffrey Youett, Elena Rysikova, \\
                                   & Philipp Droll, Yasemin Günther (später), Melf Johannsen,\\
                                   & Tahnee Herzig, Mario Moldenhauer, Freya Schmitz, \\
                                   & Jim Martens, Svenja Horn, Mirzo Ulugbek Khatamov \\
                                   & Martin Sievert, Armin Günther \\
             BGZM-Liste (2 Sitze) : & Jakob Pape, Karen Martirosian \\
             WiWi (2 Sitze): & Claas-Friso Hente, Kay Zöllmer \\
             Unicorns (5 Sitze): & Katharina Kucza, Johannes Peplow, Annkathrin Löffler, \\
                                 & Andreas Hartkamp, Marielle Hermstrüwer \\
             Liste LINKS (3 Sitze): & Gunhild Berdal, Till Petersen, Sinah Mielich \\
             HWP (2 Sitze): & Samet Gunay, Ajdina Karahasan \\
             MIN (4 Sitze): &  \\
             SDS* (3 Sitze): & Mena Winkler, Jacob Petersein, Artur Brückmann \\
             Bart-LISTE (2 Sitze): & Dominic Laumer \\
             LHG (1 Sitz): & Tobias Heisig \\
             harte zeiten (1 Sitz): & Tobias Berking \\
             Jura (1 Sitz): &  \\
             AL (2 Sitze): & Karima Schulze, Henri Weber \\
            & \\
        Entschuldigt: & Lotte Rullkötter (MIN), Timo Zeimet (Bart), Thea Wahlers (MIN)\\
                      & Jan Detampel (MIN), Ailina Salten (MIN)\\
                      & \\
        Unentschuldigt abwesend: &  Johann Baumhoefener (Jura)\\
                                &\\
        Rücktritte: & \\
    \end{tabular}
    \newpage
    %%% Vorgeschlagene Tagesordnung %%%
    \underline{Vorgeschlagene Tagesordnung}
    \begin{enumerate}[label={\textbf{Top \theenumi}},leftmargin=*]
        \setcounter{enumi}{10}
        \item Bestätigung des RBCS-Referenten (V1617-010)
        \item \#uhh hilft (V1617-007,V1617-012)
        \item
            \begin{enumerate}
                \item Rechenschaftsbericht des amtierenden AStA
                \item Fragen und Diskussion
                \item Entlastung des AStA
            \end{enumerate}
        \item Wahl des neuen AStA-Vorstandes
            \begin{enumerate}
                \item Diskussion VS-Thesen (V1617-011)
                \item Wahl des AStA-Vorstandes
            \end{enumerate}
        \item Bestätigung der AStA-Referent*innen
        \item Antrag Ältestenrat zur Beauftragung des SWOGA (V1617-013)
        \item Verschiedenes
    \end{enumerate}

    \newpage

    {\Large \underline{Teil A}}

    \subsection{Formalia}

    \subsubsection{Geschäftsbericht des Präsidiums}

    Ramon berichtet, dass das Präsidium sich als gleichberechtigt wahrnimmt und
    daher die Posten Präsident*In, Vizepräsident*In und Schriftführer*In nicht
    verteilt hat.

    Ferner erklärt Ramon, dass das Präsidium in Bezug auf den Verlust des Mandates
    einen Sitzungstag und nicht etwa eine Sitzung als das wichtige
    Entscheidungskriterium ansieht.

    Ausschusskonstituierungen sind durch und Ergebnisse können auf der StuPa-Website
    eingesehen werden. Sofern die Zusammenarbeit mit den Ausschüssen funktioniert,
    werden auch in Zukunft die Protokolle auf der Website erscheinen.

    \subsubsection{Feststellung der endgültigen Fassung des Teils B der Tagesordnung}

    Ramon schlägt vor, dass die Bestätigung des RBCS-Referenten dann stattfindet,
    wenn dieser anwesend ist und der dann stattfindene TOP fertig ist.

    Till fragt nach, ob TOP 13 als einheitlicher TOP gesehen wird. Ramon antwortet,
    dass das Präsidium dies so sieht.

    Danach kein Widerspruch zum Vorschlag.

    %%% BESCHLOSSENE TAGESORDNUNG %%%
    \underline{Endgültige Fassung der Tagesordnung}
    \begin{enumerate}[label={\textbf{Top \theenumi}},leftmargin=*]
        \setcounter{enumi}{10}
        \item Bestätigung des RBCS-Referenten (V1617-010)
        \item \#uhh hilft (V1617-007,V1617-012)
        \item
            \begin{enumerate}
                \item Rechenschaftsbericht des amtierenden AStA
                \item Fragen und Diskussion
                \item Entlastung des AStA
            \end{enumerate}
        \item Wahl des neuen AStA-Vorstandes
            \begin{enumerate}
                \item Diskussion VS-Thesen (V1617-011)
                \item Wahl des AStA-Vorstandes
            \end{enumerate}
        \item Bestätigung der AStA-Referent*innen
        \item Antrag Ältestenrat zur Beauftragung des SWOGA (V1617-013)
        \item Verschiedenes
    \end{enumerate}

    \newpage
    \setcounter{subsection}{10}
    \vspace{1cm}
    {\Large \underline{Teil B}}

    \subsection{Bestätigung des RBCS-Referenten}

    Dieser TOP wurde, wie bei der Feststellung der Tagesordnung beschlossen,
    aufgerufen nachdem Lars anwesend und der dann behandelte TOP fertig war.
    Konkret war dies der Fall nach TOP 13.

    \subsubsection{Bestätigung der Wahlniederschrift}

    Geoffrey stellt die Wahlniederschrift vor. Es gab leider eine geringe
    Wahlbeteiligung.

    Keine Fragen zum Protokoll.

    Die Wahlniederschrift wird einstimmig angenommen.

    \subsubsection{Bestätigung des Referenten}

    Lars Dierksen stellt sich vor. Wahlbeteiligung ist seit längerer Zeit gering.
    Problem ist, dass die Gruppe sehr divers und schwer zu zentralisieren
    ist.

    Beratung wird stark frequentiert. Lars hat im vergangenen Semester
    einen Schwerpunkt auf Beratung gelegt. Bisher zwei Beraterinnen,
    die auf eigenen Wunsch ehrenamtlich arbeiten. RBCS versucht mehr
    Projektkräfte zu bekommen, um mehr Veranstaltungen anbieten zu können.

    Ergonomiekurs war ausgebucht. Selbstverteidigung für Frauen war jedoch
    von nur 3 Teilnehmerinnen besucht.

    Dieses Semester ist die größte Veranstaltung der Berater*Innen-Workshop
    für ungefähr 10 Studierende. Workshop findet an einem Wochenende
    Anfang Juni statt. Nächsten Monat findet auch das Ergonomieseminar noch einmal
    statt.

    Till ist der Meinung, dass die Studierenden die Arbeit des Referats gut finden,
    aber keine Notwendigkeit sehen sich selbst zu engagieren. Die Konzentration
    auf Beratung und "Stellvertreterpolitik" könnte da die Ursache sein. Er fragt,
    ob das Referat vor hat in Zukunft mehr politische Arbeit zu machen.

    Artur sagt, dass aktuelle politisch-kulturelle Situation dazu führt, dass
    sich Leute ungern der Gruppe der Behinderten und chronisch-kranken
    Studierenden (BCS) zuordnen.

    Golnar fragt, ob es sinnvoll wäre sich mit der Geschichte in dem Bereich
    BCS auseinanderzusetzen.

    Lars sagt, dass er es ähnlich sieht, wie Till, Artur und Golnar. Das RBCS
    setze sich durchaus auch politisch ein. Mehr politische Arbeit sei
    geplant, aber es sei immer schwer Leute zu erreichen.

    \textit{Zeitansage: 21.00 Uhr}

    \textit{Lars wird einstimmig als Referent bestätigt.}

    \subsection{\#uhh hilft}

    Golnar stellt den Antrag V1617-007 zu \#uhh hilft vor.

    Jennifer stellt den Ersetzungsantrag V1617-012 zu \#uhh hilft vor.

    Geoffrey fragt aus welcher Motivation heraus der Ersetzungsantrag kommt.
    Ihm fehlt die Situationsanalyse in dem Ersetzungsantrag und fragt die
    Antragssteller*Innen, ob dies beabsichtigt ist und bittet sie generell
    darauf einzugehen.

    Till erläutert, warum die Situationsanalyse und Begründung wichtig sind und
    ruft dabei Punkte aus der Situationsanalyse und Begründung auf.

    Jennifer erklärt, dass sie das Ergebnis des Antrags gut finden, aber nicht
    die Annahmen und deswegen haben sie einen Ersetzungsantrag gemacht.

    Golnar sagt, dass sich dann die Frage stellt, warum das Ergebnis gut gefunden
    wird, wenn die Annahmen nicht geteilt werden. Sie erläutert, was die inhaltliche
    Ausrichtung der AG Flüchtlinge ist. Sie sagt, dass eine politische Begründung
    des Antrags notwendig ist, um das Programm auch durchzusetzen.

    Gunhild findet es bemerkenswert, dass die Uni gegen das Prinzip des teile
    und herrsche vorgeht. Sie findet es verlogen zu sagen, dass man die Ziele
    des Antrags unterstützt, aber nicht bereit ist die Politik der Bundesregierung
    inhaltlich zu kritisieren. Sie führt fort, dass DaF eine historische Bedeutung
    an der Uni hat und wieder in der Germanistik angesiedelt werden sollte.
    Das Anfängerniveau sollte massiv ausgebaut werden. Gunhild findet, dass dieser
    Ausbau von DaF allen zugute kommt - gerade vor dem Hintergrund, dass häufig
    Flüchtlinge als Quell vielen Übels gesehen werden.

    Franziska findet es bemerkenswert, dass das Programm (\#uhh Hilft) die Flüchtlinge
    in alle Facetten der Universität aufnimmt. Sie hat das Programm so verstanden,
    dass es u.a. beabsichtigt Flüchtlingen einen Überblick über die politische
    Situation in Deutschland zu verschaffen. Sie meint, dass der Ausblick, dass dieses
    Programm nur den Anfang einer Bewegung darstellen kann, in dem Ersetzungsantrag
    fehlt. Man habe eine Position erreicht, in der dieses Programm richtig und gut
    sei, die Welt aber noch nicht.

    Golnar hat einen Änderungsantrag eingereicht, welcher von den Antragsstellern
    übernommen wird. Der Antrag besagt, dass nach Zeile 36 folgendes eingefügt
    werden soll:

    \blockquote{Darüber hinaus fordert das Stupa die Bürgerschaft auf, den Hamburger
    Hochschulen für den bedarfsgerechten Ausbau von DaF und die Durchführung
    der Studien(vorbereitungs)programme für Geflüchtete
    dauerhaft bedarfsgerechte Mittel zusätzlich zuzuweisen.}

    Es kommt zur Abstimmung. Dabei wird der Ersetzungsantrag zuerst abgestimmt,
    da dieser der weitgehendste Änderungsantrag ist.

    \textit{Mit 6:Mehrheit:1 ist der Ersetzungsantrag nicht angenommen.}

    \textit{Es kommt zur Abstimmung des Gesamtantrags. Bei 2 Enthaltungen wird
    der Antrag einstimmig angenommen.}

    \textit{Zeitansage: 18.57 Uhr}

    \textit{Ramon stellt fest, dass Lars noch nicht anwesend ist.}

    \subsection{Entlastung des AStA}

    \subsubsection{Rechenschaftsbericht des AStA}

    Moritz hält seinen Rechenschaftsbericht. Er "warnt" davor, dass dieser
    Dinge aus 2 Legislaturperioden enthalten könnte.
    Moritz erläutert, wie die Rechenschaftsberichte aufgeteilt sind. Seiner
    ist hauptsächlich aus der Vorstandsperspektive heraus verfasst.

    Anschließend stellt Moritz die Gliederung seines Rechenschaftsberichts
    vor. Dieser ist in vier Punkte unterteilt. Ergänzend möchte er am Ende noch
    etwas zu Baufragen sagen.

    \begin{enumerate}
        \item Formale Aufgaben
        \item inhaltliche Richtung andiskutieren und vorbereiten (politischer
        Vorstand)
        \item Repräsentation des AStA gegenüber der Uni
        \item Repräsentation des AStA gegenüber der Stadt
    \end{enumerate}

    \textbf{1. Formale Aufgaben}

    Der Vorstand verlässt sich auf die Festangestellten. Daher muss der Vorstand
    eigentlich nur Sonderfälle behandeln. Regelfälle werden "quasi wie von selbst"
    durch die Festangestellten übernommen.

    Es bestand die Zielsetzung die Urlaube so zu legen, dass immer mindestens ein
    Zeichnungsberechtigter anwesend war.

    Für Rechtsstreitigkeiten musste viel Zeit aufgewendet werden. Vor diesem
    Hintergrund hat Moritz einige Überlegungen angestellt:

    \begin{enumerate}
        \item Eine der Festangestellten (z.B. Maike) sollte sich mit
        Arbeitsrecht beschäftigen. Ansonsten ist es schwer im AStA an Kompetenz
        für Arbeitsrecht zu kommen.
        \item Der AStA sollte öfter darüber diskutieren, warum die VS etwas tun
        sollte. Danach ist es häufig einfacher zu sagen, warum die VS etwas tun
        darf. Diese Überlegung steht vor dem Hintergrund, dass es häufig
        Rechtsstreitigkeiten über das politische Mandat der Verfassten Studierendenschaft
        gibt.
        \item Erfahrungen weitergeben.
    \end{enumerate}

    \textbf{2. inhaltliche Richtung andiskutieren und vorbereiten (politischer
    Vorstand)}

    Der Vorstand hätte Konflikte mehr im Vornherein besprechen müssen.
    Der politische Vorstand sollte in Zukunft stärker genutzt werden.

    \textbf{3. Repräsentation des AStA gegenüber der Uni}

    Die Repräsentation war einfacher, da man sich nicht mehr vorstellen musste.
    Der AStA ist ein wichtiger Gesprächspartner in der Uni. Der Kontakt zu den
    Lehrenden sollte weiter ausgebaut werden. Es wurde Anerkennung erworben und
    der AStA wird geschätzt. Selbst Leute, die dem AStA inhaltlich nicht
    zustimmen, freuen sich über eine aktive VS.

    \textbf{4. Repräsentation des AStA gegenüber der Stadt}

    Mit NOlympia wurde ein großer Eindruck hinterlassen. Außerdem waren
    Schulklassen auf dem Campus und der AStA war auch in Schulveranstaltungen.
    Eigene inhaltliche Veranstaltungen sind aber noch ausbaufähig.
    Als positives Beispiel von eigenen Veranstaltungen wird "Kapitalismus und
    Nachhaltigkeit" genannt.

    \textbf{Bau}

    Wie angekündigt berichtet Moritz am Ende noch über Baufragen.
    Dem AStA war es wichtig die Nutzerinteressen in dem Verfahren zu verankern
    und alle Gruppen an entsprechenden Gesprächen zu beteiligen.
    Dabei ist die Bündelung der Interessen und eine Vermittlung ebenfalls wichtig.
    Moritz hat dieser Part am meisten Spaß gemacht.

    \hrule

    Vincent fährt fort mit seinem Rechenschaftsbericht.
    Der Vorstand vertritt den AStA nach innen und außen. Er hat formale und
    politische Verantwortung. Vincent sieht den Finanzvorstand als Kristallationspunkt
    zwischen formalen Vorgaben und inhaltlicher Ausrichtung. Vincent hat viel
    zu Arbeitsrecht und Finanzrichtlinien gelernt.

    Linker AStA eckt auch an, aber das ist wichtig. Bürgerschaftsanfragen und
    Klagen sind nicht immer spaßig, aber teilweise notwendig. Dabei ist es wichtig
    relevante Rechtsangelegenheiten von unwichtigen zu unterscheiden, um die
    relevanten Angelegenheiten führen zu können.

    Vincent listet einige Punkte der inhaltlichen Arbeit des AStA auf. Der Vorstand
    hat in vielen Bündnissen mitgearbeitet. Eine engere Zusammenarbeit mit dem fzs
    sollte angestrebt werden.

    Das Amt als Vorsitzender ist sehr vereinnahmend. Vincent ist seinem Anspruch
    eine Balance zwischen formalen Vorgaben und egalitärer Arbeitsteilung zu finden,
    nicht immer gerecht geworden.

    \hrule

    Es folgt der Rechenschaftsbericht des politischen Vorstands, welcher durch
    Franziska gehalten wird. Sie sagt, dass im vergangenen Jahr vieles
    erreicht wurde. Beispielhaft dafür wird der Dies Academicus erwähnt.
    Als Leitspruch ihres Rechenschaftsberichts lässt sich festhalten,
    dass AStA Arbeit sich immer wie eine Dies Academicus Vorbereitung anfühlen
    und jeder Tag an der Universität wie ein Dies Academicus sein sollte.

    \hrule

    Timo hält seinen Rechenschaftsbericht als Finanzreferent. Er wurde erst im
    Januar Referent und sein Bericht kann dadurch nur die letzten 3 Monate umfassen.
    Er ruft noch einmal den Anspruch auf, mit dem er antrat: "so inhaltlich wie
    möglich, so formal wie nötig". Anfangs hatte er nur theoretisches Wissen
    über die Vorgänge im Finanzreferat. Durch eine Große Anfrage der FDP hat
    er viel gelernt.

    Timo ist sich nicht sicher, ob das Finanzreferat das kommunikativste Referat
    ist, allerdings hat er jeden mal gesehen und gesprochen.

    Es folgt eine kurze Darstellung der finanzielle Situation. Alle Töpfe sind im
    Plus, außer der Topf des Präsidiums.

    Er hat nicht das Gefühl, dass Leute auf der Strecke bleiben wegen Bürokratie.

    \hrule

    Es folgt der Rechenschaftsbericht des HoPo-Referats. Dieser wird von Philipp
    und Mena gehalten.

    Philipp sagt, dass sich aufgrund der langen Zusammenarbeit viel Handlungsfähigkeit
    erarbeitet wurde. NOlympia sei dafür ein gutes Beispiel. Es gibt gute
    Möglichkeiten der Zusammenarbeit mit anderen Studierendenschaften in Hamburg
    (z.B. über die LandesAStenKonferenz). Diese gute Zusammenarbeit sollte auf
    die fzs übertragen werden - gerade vor dem Hintergrund, dass die VS der Uni
    Hamburg im bundesweiten Vergleich sehr viel Möglichkeiten hat.

    \textit{Zeitansage: 19:30}

    Mena meint, dass vieles bereits gesagt wurde. Sie findet die Ausrede "ich
    muss noch lernen" sehr schlecht. Lernen sollte etwas Positives sein. Die Ausrede
    zeuge von einem degenerierten Verständnis von Lernen. Sie erwähnt in dem
    Zusammenhang die Rahmenprüfungsordnung (RPO), welche letztes Jahr entwickelt wurde
    und sich jetzt im Endprozess befindet. Mit der RPO sei eine Verbesserung der
    Studienbedingungen möglich.

    Ferner sieht sie ein Verbesserungspotenzial im AStA bzgl. der grundlegenden
    inhaltlichen Ausrichtung. In Zukunft sollte gemeinsames Lehren und Lernen
    weiter angegangen werden.

    \hrule

    Es folgt der Rechenschaftsbericht von Artur aus dem Referat für Soziales.

    Zu Beginn liest Artur einen Artikel von der ZEIT zu Olympia vor und bewertet
    diesen anschließlich. Außerdem geht er noch einmal auf die Arbeit des AStA
    zu Olympia ein. In Zukunft sollten wissenschaftlich-aufklärerische Argumente
    stärker in die Gesellschaft getragen werden.

    Artur zitiert die AStA-Leitlinien von 2015 und geht inhaltlich auf die
    genannten Punkte ein.

    Er sieht Schwierigkeiten dabei das Verhältnis zwischen Hilfe im Bestehenden
    und strukturelle Verbesserung gut hinzubekommen. Zudem hätten
    Auseinandersetzungen z.B. mit dem Studierendenwerk zu häufig im Hinterzimmer
    stattgefunden.

    Insgesamt müsse der AStA stärker herauskommen und inhaltliche Arbeit
    stärker verallgemeinern.

    \hrule

    Mario folgt mit dem Rechenschaftsbericht für das Referat für Ökologie
    und Nachhaltige Entwicklung.

    Mario hat im Oktober als Quereinsteiger im Referat für Ökologie und
    Nachhaltige Entwicklung begonnen. Anfangs musste er sich erst einmal
    in die inhaltliche Arbeit des Referats einarbeiten und dankt allen,
    die ihm dabei geholfen haben.

    Er hat zusammen mit dem restlichen Referat die bereits geplanten Veranstaltungen
    umgesetzt. Neue Veranstaltungen konnten leider nicht im angedachten
    Umfang durchgeführt werden.

    In Bezug auf seinen Anspruch zieht er ein gemischtes Fazit. Mario findet
    die anstehende Zusammenlegung des Referates sinnvoll.

    \hrule

    Abschließend folgt der Rechenschaftsbericht von Stefan Seld.
    Unter anderem geht Seld dabei auf ein Radio ein, bei dem Themen des AStA
    diskutiert werden. Dieses Radio wird Freies Senderkombinat genannt.

    \hrule

    Ramon schlägt vor jetzt die Beschlussfähigkeit festzustellen und anschließend
    eine Präsidiumspause zu machen. Während der Pause sollen die Rechenschaftsberichte
    von Till Mehrpohl und Carina gedruckt werden, sodass das Parlament diese
    zur Kenntnis nehmen kann.

    Till sagt, dass es sinnvoll wäre die Berichte jetzt schon zu drucken, um sie
    dann in der Pause lesen zu können.

    Moritz schlägt vor, dass er die Berichte ausdrucken kann, da er für die
    Beschlussfähigkeit nicht erheblich ist.

    Da einer der Berichte per LibreOffice nicht gelesen werden kann,
    wird technisch bedingt dem Vorschlag Ramons gefolgt.

    \textit{Beschlussfähigkeit mit 35 von 47 Parlamentarier*Innen gegeben.}

    \textit{Präsidiumspause bis 20.20 Uhr}

    Es ergibt sich in der Pause, dass Carinas Rechenschaftsbericht
    auch per Word nicht geöffnet werden kann.

    Oliver Vornfeld liest den Rechenschaftsbericht von Till Mehrpohl vor.
    Der Bericht kann dem Anhang entnommen werden.

    \subsubsection{Fragen und Diskussion}

    Till stellt eine Frage an Moritz und Vincent: An welchen Stellen habt
    ihr den Widerspruch zwischen formaler Hierarchie und egalitärer
    Arbeitsweise gesehen?

    an Moritz: Wie kann die "Kür" der inhaltlichen Arbeit stärker in der
    Arbeit des Vorstandes stattfinden?

    an Moritz: Ziehst du die Konsequenz die FSRK stärker in den AStA einzubinden?

    an Vincent: Wieso siehst du einen Widerspruch zwischen "Reformismus"
    und "Aktionismus"?

    an Timo: Hast du eine Vorstellung, wie aus dem Umstand des Überblicks
    über alle Referate, etwas Sinnvolles gemacht werden kann?

    Moritz beantwortet die ersten beiden Fragen. Die Aufgabe des Vorstandes
    ist es an allen Stellen mitzuarbeiten. Service ist nur solange Sklaverei,
    solange man nicht an inhaltlicher Arbeit beteiligt ist. Strukturelle
    Arbeit wurde im Vorstand gebündelt, weil Arbeitsteilung manchmal auch
    sinnvoll ist.

    Moritz beantwortet die Frage zur FSRK und sagt, dass die Beteiligung der FSRe
    in Hamburg im Vergleich zu anderen Städten nicht optimal gelöst ist. Eine
    Lösung könne jedoch nicht an der FSRK festgemacht werden. Die ideale Lösung
    sei noch unklar, es bestehe aber klarer Handlungsbedarf.

    \textit{Zeitansage: 20.30 Uhr}

    Anschließend beantwortet Vincent die Frage zum Verhältnis damit, dass er
    zu Beginn seiner Vorstandsarbeit viel zum Verhältnis gelernt habe.
    Gerade das Verhältnis zu den Festangestellten sei dabei zu nennen.
    Außerdem müsse zwischen Amt und Person unterschieden werden. Der Vorstand als
    Amt ist Arbeitgeber der Festangestellten -- nicht die Personen im Vorstand.

    Vincent beantwortet die Frage zum Widerspruch zwischen "Reformismus"
    und "Aktionismus".

    Timo beantwortet die Frage zur Übersicht, die man als Finanzreferent*In
    hat.

    Gunhild fragt, wo die Gewichtung liegt (Anerkennung oder Anecken).

    Till sagt, dass das Finanzreferat der Ort sei, wo potentielle Konflikte
    oder Herausforderungen rechtzeitig thematisiert werden könnten.
    Bei allem, was man tut, ginge es um Menschen. Gremien seien daher kein
    Selbstzweck, sondern für Menschen da.

    Moritz kann die Frage von Gunhild nicht beantworten, da er den Widerspruch
    nicht sieht. Er ist der Meinung, dass man sich Anerkennung durch ein sinnvolles
    Maß an Anecken erarbeite.
    Beispielhaft führt er eine Diskussion um eine Ordnung an, wo es keinen Sinn
    mache pauschal gegen eine Ordnung zu sein.

    Vincent meint, dass es sinnvoll sei immer eine Gesprächsgrundlage mit
    Leuten zu haben.

    \textit{Zeitansage: 20.46 Uhr}

    \subsubsection{Entlastung des AStA}

    Bei einigen Gegenstimmen und 3 Enthaltungen wird der AStA
    mehrheitlich entlastet.

    \subsection{Wahl des neuen AStA-Vorstandes}

    \subsubsection{Diskussion VS-Thesen}

    Franziska macht einen Aufschlag zur Diskussion und stellt
    den Antrag vor.

    Philipp geht auf die Bemühungen um Kapazitätsrecht und Exzellenzinitiative,
    sowie die historische Entwicklung um Wissenschaftsauffassung ein.

    Es soll unter TOP 15 keine einzelnen Vorstellungen der Referent*Innen geben.

    \textit{Fraktionspause vom RCDS im Umfang von 10 Minuten bis 21.25 Uhr.
    Pause wird durch dynamisch genommene Präsidiumspause bis 21.35 Uhr
    verlängert.}

    Mena, Geoffrey und Janis gehen nach vorne. Mena sagt, dass wir uns in einer
    Entscheidungssituation befänden. Diese Situation solle sich auch auf Aktivitäten
    im Hochschulumfeld auswirken. Dies sei im VS-Thesentext so beschrieben.

    Es wird erklärt, dass HoPo, Kultur und Öffentlichkeitsarbeit thematisch
    zusammengelegt werden. Die Universität sei als Lebensraum zu betrachten.
    Dementsprechend sei sie ein Lebensmittelpunkt, in dem man sich begegnete,
    sowie sich und letztendlich auch die Welt verändere.

    Janis sagt etwas zum kulturpolitischen Bereich. Wichtig dabei sei, dass sich
    der AStA nicht als Dienstleister sondern als Impulsgeber verstehe.
    Der AStA könne insofern als Katalysator gesehen werden und die zentrale
    Stelle nutzen, um alle Bestrebungen zu fördern.

    Geoffrey stellt die Öffentlichkeitsarbeit vor. Ein Paradebeispiel für sein
    Konzept sei die Holzwirtschaft, die davor steht abgewickelt zu werden. Es
    müsse daraufhin gewirkt werden, dass das Silodenken ("mein Studium, meine Fakultät")
    überwunden wird. In dem Sinne erinnert Geoffrey an die Demo am 3.5., wo es
    nicht nur um die Holzwirtschaft, sondern um die Finanzlage an der Uni insgesamt
    geht. Er schließt damit, dass eine gemeinsame Position nicht durch eine
    Website entstehe, sondern durch das gemeinsame Gespräch.

    Als nächstes hält Seld einen Redebeitrag. Er hebt hervor, dass es in dem
    Thesenpapier um Solidarität gehe. Er liest in dem Papier auch, dass sich die
    verschiedenen Kämpfe miteinander verbinden lassen.

    \textit{Zeitansage: 21.45 Uhr}

    Seld sagt, dass wir in einer Umwelt leben, die stark von Menschen beeinflusst
    wird. Er möchte eine Kampagne anstoßen, um den HVV kostenlos zu machen.

    Franziska hält einen Redebeitrag. Unter anderem sagt sie, dass mit dem Text
    eine wichtige Analyse getroffen worden sei, die auch für die zukünftige
    Arbeit Relevanz habe.

    Jacob Petersein geht auf den "antifaschistischen Impetus" des Textes
    ein.

    Till geht auf den letzten Absatz des Textes ein. Er stellt die Radikalität
    des Vorsatzes dar, die Solidarität auszubauen. Er geht auf zwei Personen
    ein, die früher Teil eines ausgesprochen antilinken AStAs waren, aber erkannt
    hätten, dass die Arbeit in der VS extrem befreiend sei. Till findet es wichtig
    zu reflektieren, wie wichtig Anerkennung sein darf, da es den Anerkennenden
    eine Institution zuspreche. Die Alternative zur kapitalistischen Gestaltung
    der Umgebung müsse die menschenwürdige Gestaltung der Umgebung sein. Eine
    Absage zu der Dialektik (Anecken vs. Anerkennung) sorge dafür, dass
    man in eine reaktionäre Argumentation abrutschen könne. Dies sollte vom
    AStA berücksichtigt werden.

    Armin findet, dass die ökologischen Fragen in Zukunft eine größere Rolle
    einnehmen sollten.

    \textit{Zeitansage: 22.00 Uhr}

    Artur sagt, dass man sich jeglicher rechter Hetze zur Wehr setzen müsse.
    Studierende ließen sich noch zu häufig einreden, dass sie privilegiert
    seien und das Engagement in AStA und Co. nur vorübergehend und eine Phase
    sei. Er geht auf den US-Wahlkampf mit Donald Trump und Bernie Sanders ein.
    Studierende sollten es als gemeinsamen Kampf verstehen, wenn z.B. die
    Krankenschwestern von einem Krankenhaus streiken.
    Artur verweist auch noch einmal auf NOlympia und die laufenden Anstrengungen
    die Arbeiten ins Italienische zu übertragen, damit die Errungenschaften
    in Hamburg auch auf Rom übertragen werden können. Dort soll es jetzt auch
    ein Referendum zu Olympia geben.

    Franziska hält einen weiteren Redebeitrag. Sie geht auf die Beeinflussbarkeit
    der Hamburger Regierung ein und erwähnt die Enthaltung von Fegebank
    bei der Exzellenzinitiative.

    \textit{Es folgt die Abstimmung über den Text. Bei einigen Gegenstimmen wird
    der Text mehrheitlich angenommen.}

    Es wird eine persönliche Erklärung von Ramin, Jennifer, Claas-Friso
    und einigen mehr, deren Unterschrift nicht entziffert werden konnte, vorgelesen:

    \blockquote{Wir halten die Diskussion der VS-Thesen wie vorgetragen für die
    Vorstellung des jeweiligen Parteiprogramms und fordern, dass sie auch dementsprechend
    benannt werden.}

    Anschließend wird eine persönliche Erklärung von Geoffrey verlesen:
    \blockquote{Die Opposition wieder verkennt\\
    die Wichtigkeit vom Argument,\\
    denn will man wirklich überzeugen\\
    hilfts nicht sich unterm Tisch zu beugen.\\
    Denn nur durchs Reden, Diskutieren\\
    wird inhaltlich etwas passieren.

    Eine Erklärung zum Finale\\
    wird gern auch Klingelstreich genannt.\\
    Man klopft und lächelt etwas fahle\\
    doch geht die Tür auf wird gerannt.}

    \subsubsection{Wahl des AStA-Vorstandes}

    Es gibt eine Kandidaturliste:

    \begin{enumerate}
        \item Vorstand: Philipp Droll
        \item Vorstand: Franziska Hildebrandt
    \end{enumerate}

    Philipp, Franziska und Timo stellen sich vor.

    Till stellt die Frage an Philipp, warum ausgerechnet
    Philipp Vorstand macht.

    Till stellt eine Frage an Philipp und Franziska, welche Punkte sie
    von dem Rechenschaftsbericht von Vincent und Moritz übernommen haben.

    Philipp beantwortet die Frage, warum er Vorstand macht damit, dass
    er hofft, dass es noch andere Gründe gibt, als dass er bei CampusGrün ist.
    Außerdem habe er die Arbeit im AStA sehr positiv wahrgenommen. Vor dem
    Hintergrund seiner gewonnenen Erkenntnisse sieht er es daher als sinnvoll an,
    dass er in den Vorstand geht. Er möchte tendenziell der Empfehlung von Vincent
    folgen, stärker im fzs zu arbeiten.

    Franziska geht in vielen Punkten auf die Arbeit des AStA ein.

    Es gibt keine Personaldebatte.

    Die Wahl findet offen statt.

    \textit{Mit 28:11:0 ist der AStA-Vorstand gewählt. Philipp und Franziska nehmen
    die Wahl an.}

    \subsection{Bestätigung der AStA-Referent*Innen}

    Es folgt der Vorschlag des gewählten AStA-Vorstands für die zu bestätigenden
    Referent*Innen.

    \textbf{Finanzreferent:}\\
    Timo Tomanek

    \textbf{HoPo, Öffentlichkeitsarbeit, Kultur:}
    \begin{itemize}
        \item Mena Winkler
        \item Geoffrey Youett
        \item Janis Franken
        \item Kim Isenberg
    \end{itemize}

    \textbf{Soziales, Ökologie, Soziale Bewegung:}
    \begin{itemize}
        \item Artur Brückmann
        \item Armin Günther
        \item Stefan Seld
    \end{itemize}

    \textbf{Antifaschismus, Antidiskriminierung, Kritische Wissenschaft:}
    \begin{itemize}
        \item Karima Schulze
        \item Mario Moldenhauer
        \item Annkathrin Löffler
    \end{itemize}

    \textit{Zeitansage: 22.30 Uhr}

    \textit{GO-Antrag von Melf auf Verlängerung der Sitzung bis nach TOP 17.
    Kein Widerspruch.}

    Es folgt die Vorstellung der Referent*Innen. Kim Isenberg ist aus
    gesundheitlichen Gründen leider nicht anwesend. Janis stellt daher Kim vor.
    Bei seiner eigenen Vorstellung erwähnt Janis Debatten, die im Bereich der
    Kultur anstünden.

    An die Vorstellung schließt sich eine Fragerunde an.

    Till stellt die Frage, an welcher Stelle in der Referatsstruktur die
    Beschäftigung mit der Friedensentwicklung verortet sein soll.

    Till stellt die Frage, welche Bedeutung die Personen aus Antifa und Antidis
    dem sozialen Konflikt zumessen.

    Till findet, dass die Wörter "halt" und "irgendwie" dafür sprechen, dass sich
    Seld cool findet. Gibt es Punkte, die Seld wirklich inhaltlich teilt?

    Golnar hat die Frage, was die Referent*Innen unbedingt anpacken wollen.

    Die Beantwortung dieser Fragen beginnt mit Seld.

    Seld hat Probleme mit dem Thesenpapier, v.a. mit dem Duktus und der
    Entstehungsgeschichte. Dabei geht er u.a. auf die Schuldenbremse ein.

    Annkathrin geht auf Tills Frage zu Antifa und Antidis ein. Sie sieht
    den sozialen Konflikt auch als Problem an.

    \textit{Zeitansage: 23.00 Uhr}

    Artur meint, dass Frieden ein Querschnittsthema sein sollte. Der AStA sollte
    den imperialistischen Impuls zurückweisen. Artur will unbedingt ändern,
    dass Studierenden eingeredet wird, dass sie an den neoliberalen Zuständen
    Schuld seien.

    Mena sagt, dass die Friedensfrage überall aufgehoben sein sollte.
    Sie möchte daran arbeiten, dass weniger nebeneinander her studiert und
    gearbeitet wird.

    Geoffrey sagt auch, dass Frieden übergreifend bearbeitet werden muss.
    Strukturell solle die Diskussion zur Zivilklausel weiter vorangetrieben werden.
    Kein Mensch sollte unter Kriegsbedingungen leben. Die Friedensthematik
    sei nicht nur Aufgabe des AStA, sondern der gesamten Gesellschaft.
    Geoffrey möchte als erstes die Denkweise ändern, wie Referate übergreifend
    zusammenarbeiten können. Er stellt die Frage, wie man wieder eine Universität
    hinbekommt, an der man zusammenarbeitet und nicht in Scheuklappen denkt.

    Karima antwortet für das Referat und möchte die Zusammenarbeit mit den
    teilautonomen Referaten verbessern. Außerdem soll die Arbeit des AStA
    auch außerhalb der AStA-Räumlichkeiten stattfinden und in die Universität
    getragen werden. Mario sagt, dass Frieden übergreifend ist und die Zivilklausel
    in der Kritischen Wissenschaft verortet sei.

    Es schließt sich eine zweite Fragerunde an.

    Eine Studentin namens Lena fragt, inwiefern der AStA das Thema
    Feminismus vertritt.

    Golnar fragt, wie sich der AStA zur Frage der Inklusion verhält. Desweiteren
    fragt sie, wie alle das Verhältnis von StuPa zu AStA sehen. Sie fragt ferner,
    welche Debatten im Bereich der Kultur anstünden. Golnar fragt alle Referate,
    wie sich diese an der Anti-ExzellenzIni-Kampagne beteiligen möchten.

    Karima beginnt die Antwortrunde auf diese Fragen. Sie beantwortet die Frage
    zur Diskriminierung. Es kann Personen geben, die sowohl diskriminiert werden
    als auch selber diskriminieren.

    Mario erklärt das Verhältnis von StuPa zu AStA damit, dass die Referent*innen
    vom StuPa bestätigt werden und das Verhältnis besser nicht beschrieben werden
    könne.

    Seld geht auf die Frage zu Feminismus ein. Dabei erwähnt er mehrere
    Frauenbewegungen.

    Artur sagt, dass die formale Frage, wer Referent oder Referentin ist,
    nicht relevant sei, da Frauke Petri und Beatrix von Storch beides Frauen sind,
    aber sehr frauenfeindlich agierten.

    Armin und Seld stellen klar, dass es keine Referatsmeinung zu dem
    Thema gibt und es unterschiedliche Positionen im Referat dazu gibt.

    Janis beantwortet die Frage zu den anstehenden Debatten. Beispielsweise
    stelle sich die Frage, ob "tanzen gegen rechts" genug sei.

    Mena geht auf die ExzellenzIni ein.

    \textit{Zeitansage: 23.30 Uhr}

    Beispielsweise geht sie darauf ein, dass eine Person ihr gesagt hat, dass
    sie über die ExellenzIni über die Klozeitschrift des AStA herausgefunden hat.

    Die Bestätigung findet offen statt. 7 der Referent*Innen sind stimmberechtigt
    im StuPa.

    \textit{Es gibt 16 Stimmen für die Referent*Innen. Da sieben der Referent*Innen
    stimmberechtigt, aber nicht anwesend sind, werden ihre Stimmen zu ihren
    Gunsten gerechnet. Damit ergeben sich 23 Stimmen, die eine Bestätigung
    befürworten. Es gibt 6 Gegenstimmen und keine Enthaltungen.}

    Die Referent*Innen sind damit bestätigt.

    \subsection{Antrag Ältestenrat zur Beauftragung des SWOGA}

    Geoffrey erklärt den Antrag vom Ältestenrat.

    Melf sagt, dass sich der SWOGA am 11. Mai um 18:30 das nächste Mal
    trifft.

    \textit{Der Antrag wird bei einer Enthaltung einstimmig angenommen.}

    \subsection{Verschiedenes}

    Till erinnert, dass das StuPa beim letzten Mal einen Aufruf zur Teilnahme
    an der Erinnerung an das Bücherverbrennen beschlossen hat. Das findet am
    10. Mai am Kaiser-Friedrich-Ufer statt. Alle sind dazu aufgerufen daran
    teilzunehmen.

    Vincent erinnert daran, dass der AStA zusammen mit anderen Gruppen
    für den 1. Mai aufruft. Am 3. Mai findet die Demonstration gegen
    die Schließung der Holzwirtschaft statt. Sie startet um 15:30 beim
    Dammtorbahnhof.

    Artur verweist auf eine Veranstaltung von Torsten Bultmann.
    Diese findet am 10. Mai um 18 Uhr statt.

    \textit{Die Sitzung wird um 23.41 Uhr geschlossen.}

    \newpage

    \section*{Rechenschaftsbericht von Till Meerpohl}

    Vor ca. einem Jahr sind wir als AStA gewählt worden, der sich in puncto
    Öffentlichkeitsarbeit vorgenommen hatte, in Publikationen über
    Auseinandersetzungen und Geschehnisse in Stadt und Universität kritisch zu
    berichten und erfolgreich weitere Menschen an der Praxis des AStA teilhaben
    zu lassen.

    So sinnvoll dieser Anspruch auch sein mag, glaube ich, dass wir eine ganz
    andere Herausforderung unterschätzt haben, aber aus der wir bestimmt einiges
    mitnehmen und lernen können. Das Referat für Öffentlichkeitsarbeit hat sich
    im letzten Jahr allerlei technischen Herausforderungen stellen müssen, bei
    denen der gestalterische Anteil manchmal etwas kurz gekommen ist.

    Die Herausforderung, dass es auf der einen Seite notwendige Arbeit gibt, wie
    das Layouten von Publikationen, die Pflege von Homepage und Facebook sowie das
    Verteilen der Publikationen und auf der anderen Seite den Anspruch eigene
    Ideen und ein Verständnis von Gesellschaft und Kritik an selbiger zu entwickeln
    und nach außen zu tragen sowie gestalterisch tätig zu sein, wurde häufig
    zugunsten ersterer beantwortet.

    So haben wir einen, wie ich finde, sehr guten Kalendasta, der sowohl in seiner
    optischen Präsenz als auch inhaltlich sehr gut war, konzipiert, uns immer
    wieder Gedanken zu der Strukturierung neuer Stiller Wörtchen gemacht und uns
    zum Relaunch der Homepage lange beraten, wie wir das, was der AStA tut und
    vertritt, besser im Webauftritt verdeutlichen können. Auch wenn das nicht
    immer nur etwas mit Gesellschaftskritik zu tun hatte, waren das Situationen
    in denen wir gemeinsam Dinge entwickelt haben und sich nicht das Gefühl von
    dem Abarbeiten von Aufgaben zur Erhaltung des laufenden AStA-Betriebs einstellte.

    Ich glaube, dass die Klärung des Umgangs mit solchen Aufgaben, die zweifelsfrei
    notwendig sind und getan werden müssen, die Arbeit im Referat für
    Öffentlichkeitsarbeit beleben könnte.

    Auch wenn das Augenmerk gerade mehr auf dem zu verbessernden lag, möchte ich
    doch noch einmal ganz besonders Oli und Wiebke für eine gute und schöne Zeit
    danken. Ich habe den Umgang mit euch immer sehr geschätzt und glaube auch
    viel von euch gelernt zu haben.

\end{document}
